% TOP_PROBLEM: {"problemId": "problem.toms-winter-strict-comparison-implies-z-stability", "canonicalTitle": "Toms–Winter: Strict Comparison Implies Z-Stability", "releaseRank": 380, "primaryDomain": "algebra_representation_category", "primaryDomainLabel": "Algebra, representation & category theory", "displayStatus": "open_with_solved_subcases", "releaseStatus": "open_with_solved_subcases"}
% !TeX program = lualatex
% ATTEMPT_STATUS: unresolved
% Model: GPT-6 Astra Ultra; continuation dated 27 September 2026.
\documentclass[11pt]{article}
\usepackage[margin=25mm]{geometry}
\usepackage{fontspec}
\setmainfont{Latin Modern Roman}
\usepackage{amsmath,amssymb,mathtools}
\usepackage{unicode-math}
\setmathfont{Latin Modern Math}
\usepackage{xurl,hyperref}
\hypersetup{colorlinks=true,urlcolor=blue}
\setcounter{secnumdepth}{0}
\setlength{\parindent}{0pt}
\setlength{\parskip}{5pt}
\setlength{\emergencystretch}{3em}
\begin{document}
\section*{\texorpdfstring{Toms–Winter: Strict Comparison Implies Z-Stability}{Toms–Winter: Strict Comparison Implies Z-Stability}}\label{380}
\subsection{Definitions and mathematical statement}
Let $A$ be a separable simple nuclear C*-algebra, not isomorphic to an algebra of compact operators; no unit is required. Its Cuntz semigroup $\operatorname{Cu}(A)$ consists of positive elements of $A\otimes\mathcal K$ modulo mutual Cuntz subequivalence, where $a\precsim b$ means $x_j^*bx_j\to a$ for some sequence. In the selected strict-comparison convention, $x\le y$ whenever every additive order-preserving functional $\varphi:\operatorname{Cu}(A)\to[0,\infty]$, preserving zero and increasing sequential suprema, satisfies $\varphi(x)\le\varphi(y)$, with strict inequality whenever $0<\varphi(y)<\infty$. The conjecture asserts
\[
 A\text{ has strict comparison}\quad\Longrightarrow\quad
 A\cong A\otimes_{\min}\mathcal Z,
\]
where $\mathcal Z$ is the Jiang--Su algebra. Separability means a countable norm-dense subset; simplicity means no nontrivial closed two-sided ideals; nuclearity is the completely positive matrix-factorization approximation property. No UCT, finite nuclear dimension, or extra trace-space hypothesis is appended.
\par\smallskip\textit{Formulation and concept references:} \href{https://arxiv.org/pdf/1912.04207}{[S1]}.

\subsection{Short English statement}
Does strict comparison force every separable simple nuclear non-elementary C*-algebra to absorb the Jiang--Su algebra, including nonunital cases?
\subsection{Research attempt}
\paragraph{The remaining implication and its quantifiers.}
The exact comparison convention in the question is the one printed
after Conjecture 5.1 of [S1]. Its functionals can take the value
$+\infty$; the strict inequality is required only at finite,
strictly positive values of the comparison target. Replacing this
by a condition on a single normalized trace would change the
problem, especially in the nonunital setting.

The inspected Theorem 5.6 of [S1] proves the implication in the
unital tracial case with the additional hypothesis of uniform
property $\Gamma$. That additional hypothesis is not part of
this catalogue entry. The argument below gives explicit comparison
calculations, constructs uniform central trace divisions in an
infinite tensor-tail subcase, and identifies the missing passage
from order information to approximately central structure. It
does not assume finite nuclear dimension, UCT or a regular trace
boundary in order to declare the target solved.

\paragraph{What Cuntz comparison means in the elementary test case.}
For positive compact operators $a,b$ on a separable infinite
Hilbert space, Cuntz subequivalence is exactly support-rank
comparison:
\[
 a\precsim b\quad\Longleftrightarrow\quad
      \operatorname{rank}a\leq\operatorname{rank}b,
 \tag{380.1}
\]
where infinite rank is allowed. Here rank means the dimension of
the closure of the range, not algebraic rank on a dense subspace.

To prove sufficiency, truncate the spectral decomposition of $a$
to finitely many eigenvalues. If the truncation has rank $r$,
choose $r$ positive eigenvectors of $b$, map the corresponding
eigenvectors of $a$ to them and rescale by square roots of the
eigenvalue ratios. This finite-rank operator $x$ satisfies
$x^*bx=a_{\rm trunc}$. The spectral truncations converge to $a$
in norm, proving the subequivalence. The operators $x$ may have
large norm; the definition does not require a uniform norm bound.

For necessity, if $b$ has finite rank $r$, then every $x^*bx$ has
rank at most $r$. A positive operator with at least $r+1$
positive eigenvalues cannot be a norm limit of such operators:
on the span of those eigenvectors its quadratic form is bounded
below by their positive minimum, whereas every rank-$r$ operator
has a kernel vector in that $(r+1)$-dimensional subspace. Thus
rank cannot increase in the limit. If $b$ has infinite rank,
the necessary inequality is automatic. This proves (380.1).

Since $M_d\otimes\mathcal K\cong\mathcal K$, it follows that
$\operatorname{Cu}(M_d)=\{0,1,2,\ldots,\infty\}$ with the
usual order and addition. Every finite rank $r$ is a sum of
$r$ rank-one classes, and increasing ranks have supremum infinity.

\paragraph{The non-elementary exclusion prevents an actual counterexample.}
The matrix algebra $M_d$ satisfies the printed strict-comparison
condition. To check the implication, use the rank functional
$\varphi(r)=r$. If $y$ is finite and positive, the premise forces
$x<y$; if $y=0$, it forces $x=0$; if $y=\infty$, every $x$ is
already at most $y$. Thus the required order conclusion holds
in every case in which the functional premise holds.

Nevertheless $M_d$ does not absorb the infinite-dimensional
Jiang--Su algebra: $M_d\otimes\mathcal Z$ is infinite-dimensional,
whereas $M_d$ is finite-dimensional. The same stabilization
calculation explains the elementary compact-operator case excluded
by the statement. These are not counterexamples under the selected
hypotheses; they show that deleting the non-elementary exclusion
would make the implication false.

The test also warns against equating comparison with exact
divisibility. In $\operatorname{Cu}(M_d)$ the rank-one class
cannot be twice a nonzero class. Even the weaker request
$kz\leq1\leq(k+1)z$ has no solution for $k\geq2$: the left
inequality forces $z=0$, contradicting the right. Order detected
by functionals does not automatically create intermediate
semigroup elements. Non-elementarity is necessary to exclude
this obstruction, but merely noting its necessity is not a
construction of the missing elements in a general algebra.

\paragraph{Strictly ordered projections versus central comparison.}
In the unital stably finite simple nuclear setting, normalized
trace functionals give dimension functions
$d_\tau(a)=\lim_{m\to\infty}\tau(a^{1/m})$; for a projection
$p$, this equals $\tau(p)$. The correspondence with the relevant
Cuntz functionals is part of the framework recorded in [S1].
Strict comparison then turns uniform strict dimension inequalities
into Cuntz subequivalence. Its output, however, is only a sequence
of witnesses $x_j$ with $x_j^*bx_j\to a$.

There is no approximate commutation requirement on those $x_j$.
The issue is visible in $M_3$: the rank-one projection $e_{11}$
is subequivalent to $e_{22}+e_{33}$, witnessed by $e_{21}$ in
the convention $x^*bx=a$. Yet a scalar witness cannot do this,
since $c^*(e_{22}+e_{33})c$ is supported orthogonally to $e_{11}$.
Its distance from $e_{11}$ is at least one. A witness approximately
commuting with all of $M_3$ is close to a scalar, so it cannot
make this error tend to zero.

For the last assertion, average over the finite clock-and-shift
unitaries $U^rV^s$, $0\leq r,s<3$. Their conjugation average of
an element $x$ is $\operatorname{tr}_3(x)1$. If all the corresponding
commutators are at most $\varepsilon$, this average differs from
$x$ by at most $\varepsilon$. Thus the obstruction is quantitative.
The example is elementary and excluded from the final conjecture,
but it disproves a purely formal inference from comparison
witnesses to central comparison witnesses.

\paragraph{An explicit source of central trace divisions.}
Let $U=\bigotimes_{j=1}^\infty M_2$ be the infinite tensor
product with its canonical units, and let $B$ be any unital
$C^*$-algebra. Set $A=B\otimes U$. We prove the following
approximation property directly. For a finite set $\mathcal F$
in $A$, an integer $n\geq2$ and $\varepsilon>0$, there are
orthogonal projections $e_1,\ldots,e_n$ in $A$, summing to one,
such that
\[
 \|[e_i,a]\|<\varepsilon,
 \qquad
 \left|\tau(ae_i)-\frac1n\tau(a)\right|<\varepsilon
 \tag{380.2}
\]
for every $a\in\mathcal F$, every tracial state $\tau$ on $A$
and every $i$.

Approximate $\mathcal F$ in norm within $\delta$ by elements
$a_0$ of $B\otimes\bigotimes_{j=1}^mM_2$. Take a later block
of $k$ tensor factors, isomorphic to $M_D$ with $D=2^k$.
Partition its standard basis into $n$ pieces whose sizes $r_i$
differ by at most one. The diagonal projections onto these pieces
satisfy
\[
 \sum_ie_i=1,\qquad e_ie_j=0\ (i\ne j),\qquad
 \left|\frac{r_i}D-\frac1n\right|\leq\frac1D.
\]
They commute exactly with every $a_0$.

For any tracial state on an algebra containing a commuting matrix
factor, its restriction to the first factor times that matrix
factor is the first-factor trace tensored with normalized matrix
trace. Indeed off-diagonal matrix units have trace zero after
conjugation by diagonal unitaries, and all diagonal matrix units
have equal trace after permutation conjugation. Applying this
to $a_0e_i$ gives
$\tau(a_0e_i)=(r_i/D)\tau(a_0)$, uniformly over all traces.
Consequently
\[
 \|[e_i,a]\|\leq2\delta,\qquad
 \left|\tau(ae_i)-\frac1n\tau(a)\right|
      \leq2\delta+\frac{\|a_0\|}{D}.
 \tag{380.3}
\]
Choose $\delta$ small and then $k$ large. This proves (380.2).
The rounding of ranks is handled explicitly; powers of two
need not be divisible by the prescribed integer $n$.

\paragraph{Passing this construction to a central sequence algebra.}
If $A$ is separable, choose an increasing sequence of finite sets
whose union is norm dense. Apply (380.2) with errors tending to
zero to produce sequences $e_i^{(j)}$. In the quotient
$A_\infty=\ell^\infty(A)/c_0(A)$ their classes commute with every
constant element of $A$, are orthogonal projections, and sum to
the unit. Along any ultrafilter and any sequence of tracial states,
the trace-splitting equalities pass to the limit:
\[
 \tau_\omega(ae_i)=\frac1n\tau_\omega(a)
                         \quad(a\in A).
 \tag{380.4}
\]
The construction also descends to the tracial ultrapower, since
norm-null sequences are uniformly trace-two-norm null. Thus it
supplies, in this tensor-tail setting, the uniform property
$\Gamma$ used in [S1]. The trace estimate is simultaneous over
all traces; separate constructions for each individual trace
would not have implied (380.4).

If this $A$ is simple, separable, nuclear and unital, has a
nonempty tracial state space, and has strict comparison, the
inspected Theorem 5.6 of [S1] now implies $\mathcal Z$-stability.
This is a complete conditional subcase obtained by combining an
explicit verification of the additional hypothesis with a named
external theorem. The deep passage from uniform central
divisions and comparison to absorption is attributed to that
theorem, not reproved by tensor bookkeeping.

\paragraph{Why individual traces do not give a uniform construction.}
Suppose for every trace $\tau$ one found elements splitting its
values approximately. Those elements may depend on $\tau$.
The property needed above has the order
\[
 \text{for each finite set and tolerance, choose elements that
 work for every trace.}
\]
This is different from choosing elements after fixing the trace.
Even compactness of a trace space only produces finitely many
locally valid choices; one still has to combine them inside the
algebra while preserving orthogonality, positivity and approximate
centrality. A scalar partition of unity on the trace space is
not automatically a central partition inside $A$.

Likewise, strict inequalities between dimension functions do not
by themselves give a positive numerical margin uniform over all
traces unless the relevant functions and compactness hypotheses
justify it. Dimension functions of positive elements need not
be continuous affine functions in the same way that traces of a
fixed element are. Passing from pointwise comparison to a
central construction can therefore lose two kinds of uniformity,
both across traces and across the prescribed finite set.

The infinite tensor tail avoids these difficulties for a concrete
reason: the same matrix block commutes with every approximating
head element, and every trace has exactly the same normalized
matrix-factor formula. A general nuclear approximation by
completely positive maps into matrices is not such a commuting
matrix subalgebra inside $A$.

\paragraph{Nuclear approximation is not a multiplicative splitting.}
Nuclearity provides finite-dimensional completely positive
factorizations approximating finite sets in norm. It does not
say that the matrix-factor maps are multiplicative or that their
images commute with the original finite set. Taking matrix
projections and applying a general completely positive map need
not preserve projections or orthogonality. For instance the
normalized trace map $M_2\to\mathbb C$ sends two orthogonal
rank-one projections to the same nonzero scalar $1/2$.
Thus the tail construction cannot be copied by replacing its
actual matrix block with a generic nuclear factorization.

One must obtain appropriately orthogonal order-zero or central
sequence structure with uniform tracial size. Theorem 5.6
encodes a route once uniform property $\Gamma$ is present.
Deriving that structure from only the selected comparison
hypothesis is the unproved step in this attempted approach.

\paragraph{The nonunital scope is not eliminated by adjoining a unit.}
If $A$ is nonunital and simple, its unitization contains $A$
as a proper ideal and has scalar quotient. It is generally not
simple. Applying a theorem for simple unital algebras to the
unitization therefore fails its hypotheses. Replacing $A$ by
a unital full corner also requires a suitable nonzero projection;
stably projectionless algebras do not supply one. Neither
operation is a formal way of reducing the entire printed
nonunital statement to the unital tracial theorem used above.

The functional convention in the question is designed to retain
these cases. A proof restricted to normalized tracial states
must state its restriction explicitly, rather than treating an
empty or noncompact normalized trace space as harmless.

\paragraph{Diagnostics and outcome.}
The exact checks verify the finite-rank comparison witnesses,
clock-and-shift averaging, rank partitions for tail blocks and
the error bound from rounding. They do not compute an infinite
Cuntz semigroup or certify $\mathcal Z$-absorption for arbitrary
input algebras.

We have explained why the elementary exclusion is necessary,
proved an explicit central division construction for an infinite
tensor-tail family, and used the source theorem only with its
additional unital tracial and uniform-$\Gamma$ hypotheses.
The first unresolved step for the catalogue is constructing
the requisite uniform approximately central structure from
strict comparison alone, including the separate nonunital
difficulties. The general implication remains unresolved by
this attempt.
\subsection{Sources}
\begin{itemize}
\item[S1]Castillejos, Evington, Tikuisis and White, Uniform Property Gamma, Conjecture 5.1 and following comparison definition; Schafhauser, Tikuisis and White, Nuclear C*-algebras: 99 problems, Problem XVIII.
\url{https://arxiv.org/pdf/1912.04207}.
\textit{Formulation source.}
\end{itemize}
\end{document}
